\documentclass{article}
\usepackage{listings}
\usepackage{color}

% Listings settings
\definecolor{codegray}{rgb}{0.5,0.5,0.5}
\definecolor{codepurple}{rgb}{0.58,0,0.82}
\definecolor{backcolour}{rgb}{0.95,0.95,0.92}

\lstdefinestyle{mystyle}{
    backgroundcolor=\color{backcolour},   
    commentstyle=\color{codegray},
    keywordstyle=\color{magenta},
    numberstyle=\tiny\color{codegray},
    stringstyle=\color{codepurple},
    basicstyle=\ttfamily\footnotesize,
    breakatwhitespace=false,         
    breaklines=true,                 
    captionpos=b,                    
    keepspaces=true,                 
    numbers=left,                    
    numbersep=5pt,                  
    showspaces=false,                
    showstringspaces=false,
    showtabs=false,                  
    tabsize=2
}
\lstset{style=mystyle}

\title{Monolix Installation Troubleshooting}
\author{Rohan Pandey + CHATGPT + Chen Poo}
\date{}

\begin{document}

\maketitle

\section*{Error Message 1}
When attempting to install Monolix, the following error was encountered:
\begin{verbatim}
The compiler g++ failed with the following errors:

g++ --version: Apple clang version 16.0.0 (clang-1600.0.26.4)
Target: arm64-apple-darwin24.1.0
Thread model: posix
InstalledDir: /Library/Developer/CommandLineTools/usr/bin

Please make sure that g++ command works properly and try again
\end{verbatim}

\section*{Issue 1: Clang vs GCC}
Apple's \texttt{g++} is actually a wrapper around the \textbf{Clang compiler} and not GCC, which is required by Monolix.

\subsection*{Solution Steps: Install GCC}
To install GCC via Homebrew, run:
\begin{lstlisting}[language=bash]
brew install gcc
\end{lstlisting}
This will install the GCC binaries, including \texttt{g++-14}, in \texttt{/opt/homebrew/bin}.

\subsection*{Temporarily Override g++}
To point \texttt{g++} to the newly installed GCC:
\begin{lstlisting}[language=bash]
sudo mv /Library/Developer/CommandLineTools/usr/bin/g++ /Library/Developer/CommandLineTools/usr/bin/g++_backup
sudo ln -s /opt/homebrew/bin/g++-14 /Library/Developer/CommandLineTools/usr/bin/g++
\end{lstlisting}
Verify that the correct \texttt{g++} is now in use:
\begin{lstlisting}[language=bash]
/Library/Developer/CommandLineTools/usr/bin/g++ --version
\end{lstlisting}
Expected output:
\begin{verbatim}
g++ (Homebrew GCC 14.x.x) ...
\end{verbatim}

\subsection*{Restore Clang After Installation}
After installing Monolix, restore Clang's \texttt{g++}:
\begin{lstlisting}[language=bash]
sudo rm /Library/Developer/CommandLineTools/usr/bin/g++
sudo mv /Library/Developer/CommandLineTools/usr/bin/g++_backup /Library/Developer/CommandLineTools/usr/bin/g++
\end{lstlisting}

\section*{Running the Installer in Rosetta 2}
If Monolix is not fully ARM-compatible, use Rosetta 2 for x86 emulation.

\subsection*{Install Rosetta 2}
\begin{lstlisting}[language=bash]
softwareupdate --install-rosetta
\end{lstlisting}

\subsection*{Open a Rosetta Terminal}
\begin{enumerate}
    \item Go to \textbf{Applications > Utilities} and locate the Terminal app.
    \item Right-click on Terminal, select \textbf{Get Info}, and check \textbf{Open using Rosetta}.
    \item Open Terminal.
\end{enumerate}

\subsection*{Run the Installer}
In the Rosetta-enabled terminal:
\begin{lstlisting}[language=bash]
cd /path/to/monolix/installer
open MonolixSuite2024R1-macosx-installer.app
\end{lstlisting}

\section*{Error Message 2}
During module build:
\begin{verbatim}
fatal error: 'string' file not found
#include <string>
        ^~~~~~~~
1 error generated.
\end{verbatim}

\subsection*{Solution Steps}
\begin{enumerate}
    \item Ensure you are using GCC instead of Clang:
    \begin{lstlisting}[language=bash]
    export CXX=/opt/homebrew/bin/g++-14
    export CC=/opt/homebrew/bin/gcc-14
    \end{lstlisting}
    \item Verify the standard library path. If needed, explicitly add it:
    \begin{lstlisting}[language=bash]
    export CXXFLAGS="-I/opt/homebrew/Cellar/gcc/14.0.0/include/c++/14.0.0"
    \end{lstlisting}
    \item Clean and rebuild:
    \begin{lstlisting}[language=bash]
    cd /Users/rohanpandey/lixoft/monolix/.../build
    rm -rf *
    /Applications/MonolixSuite2024R1.app/Contents/Resources/monolixSuite/bin/build
    \end{lstlisting}
    \item Run in Rosetta mode if necessary.
\end{enumerate}

\section*{Final Notes}
If you still face issues, try installing GCC 13 explicitly:
\begin{lstlisting}[language=bash]
brew install gcc@13
export CXX=/opt/homebrew/opt/gcc@13/bin/g++-13
export CC=/opt/homebrew/opt/gcc@13/bin/gcc-13
\end{lstlisting}

\newpage

\section*{Troubleshooting Steps}

When encountering the Monolix error:
\begin{verbatim}
fatal error: 'string' file not found
\end{verbatim}

follow the steps below to ensure that the correct C++ headers are accessible.

\begin{enumerate}
    \item \textbf{Install or Update Xcode Command Line Tools:}\\
    Even if you have the full Xcode IDE installed, the Command Line Tools are required. Install them by running:
    \begin{verbatim}
xcode-select --install
    \end{verbatim}

    \item \textbf{Check Your Xcode Path:}\\
    Confirm that your system is pointing to the correct Xcode version:
    \begin{verbatim}
sudo xcode-select --switch /Applications/Xcode.app/Contents/Developer
xcode-select -p
    \end{verbatim}
    The \texttt{xcode-select -p} command should return a path ending in \texttt{Developer}.

    \item \textbf{Accept Xcode License and Initialize the Toolchain:}\\
    If you recently installed or updated Xcode, you may need to accept its license and run the first launch setup:
    \begin{verbatim}
sudo xcodebuild -license
sudo xcodebuild -runFirstLaunch
    \end{verbatim}

    \item \textbf{Check macOS SDK Paths (If Needed):}\\
    Verify the availability of the SDK headers:
    \begin{verbatim}
ls /Library/Developer/CommandLineTools/SDKs
    \end{verbatim}
    If headers are missing, reinstall the Command Line Tools.

    \item \textbf{Rebuild or Relaunch Monolix:}\\
    After ensuring that your system has proper toolchain access, restart Monolix and attempt to reload the demos.
\end{enumerate}

These steps should resolve the issue where standard library headers like \texttt{<string>} are not found.

\end{document}
